\section{Montague's pragmatics and intensional logic}
\label{sec:p08-n-b-006-02}

This is the presentation of \emph{Pragmatics and Intensional Logic} in
\emph{Synthese} 22 (1970), pp.~68--94, not the same-titled
\emph{Dialectica} manifestation and not \emph{Universal Grammar}.
Montague begins with truth relative to a context, generalizes operators
to relations between indices and proposition sets, and reduces those
operators to predicates of propositions. He then introduces an
intermediate extended pragmatics and discusses the three systems'
metatheory \cite{Montague1970Pragmatics}.

The problem is indexicality: the truth of a future-tense claim about
Caesar, for example, depends on the time of utterance. Montague follows
Bar-Hillel's context-directed refinement of Morris's program for
pragmatics, but makes interpretations and context-relative truth precise.
An index need record only the relevant aspects of a context, such as a
speaker and a time. It need not be an entire model or a possible world.
The general operator construction is attributed to Montague and Charles
Howard; the reduction to predicates of propositions to Montague and
J.~A.~W. Kamp; extended pragmatics to Montague and Scott
\cite[pp.~68--69,76,82,86; note~12]{Montague1970Pragmatics}.

\subsection{Ordinary pragmatics: Definitions I--VII}

\begin{hbdefinition}{def:p08-montague-pragmatic-language}{Pragmatic language and possible interpretation}
A pragmatic language has individual variables \(v_n\), individual
constants, \(n\)-place predicate constants for \(n\geq0\), unary
operators, identity \(=\), the actuality predicate \(E\), Boolean
connectives, and individual quantifiers. No operation symbols or
descriptions occur in this initial language. Its atomic formulas are
\(E[\zeta]\), \(\zeta=\eta\), and
\(P[\zeta_0,\ldots,\zeta_{n-1}]\), with individual constants or
variables as arguments; formulas close under connectives, quantifiers,
and \(N\phi\) for an operator \(N\). We write \(\forall,\exists\)
for the source's universal and existential quantifier signs.

A possible interpretation is \(\mathfrak A=\langle A,F,R\rangle\).
Put \(I=DA=\dom A\), \(A_i=A(i)\), and
\(U=\bigcup_{i\in I}A_i\). Thus \(A_i\) is the set actual at
\(i\), whereas \(U\) is the set of possible individuals.
For an individual constant \(c\), \(F_c:I\to U\); for an
\(n\)-place predicate constant, \(F_P:I\to\mathcal P(U^n)\).
For an operator \(N\),
\[
 R_N\subseteq I\times\mathcal P(I).
\]
The source's \(\mathcal S I\) is the power set \(\mathcal P(I)\).
Constants and predicate extensions need not be actual at the index
where they are evaluated. ``The previous Pope,'' remembrance, and thinking
of an object absent from the current world illustrate the distinction.
\end{hbdefinition}

\begin{hbconvention}{def:p08-montague-empty-relation}{Relations, predicates, propositions, and empty interpretations}
A \((U_0,\ldots,U_{n-1})\)-relation is a subset of
\(\prod_{r<n}U_r\). An \((I,U_0,\ldots,U_{n-1})\)-predicate is
a function from \(I\) to such relations. ``Relation'' names an
extension, ``predicate'' an intension. Let \(\Lambda\) be the empty
sequence, identified set-theoretically with \(\varnothing\).
The empty product is \(\{\Lambda\}\); its two subsets are the truth
values \(\mathrm F=\Lambda\) and \(\mathrm T=\{\Lambda\}\).
Consequently a zero-place predicate is a proposition
\(I\to\{\mathrm F,\mathrm T\}\).

The source does not require \(U\ne\varnothing\) in a possible
interpretation. When \(U=\varnothing\), the assignment set
\(U^\omega\) is empty, and every sentence is true by the extension
definition below. Such interpretations supply no counterexamples to
sentence consequence, so restricting consequence to nonempty
interpretations gives the same relation. Evaluation at an index still
requires \(i\in I\). Statements about satisfaction by an individual
assume that individual belongs to \(U\)
\cite[p.~93, note~7]{Montague1970Pragmatics}.
\end{hbconvention}

\begin{hbdefinition}{def:p08-montague-pragmatic-extension}{Extensions, truth, satisfaction, intension, and consequence}
For \(x\in U^\omega\), set
\(\operatorname{Ext}_{i,\mathfrak A}(v_n)(x)=x_n\) and
\(\operatorname{Ext}_{i,\mathfrak A}(c)(x)=F_c(i)\).
Write \(x[n:=y]\) for coordinate replacement. Formula extensions
are subsets of \(U^\omega\), defined by
\begin{align*}
 \operatorname{Ext}_{i,\mathfrak A}(E[\zeta])
   &=\{x:\operatorname{Ext}_{i,\mathfrak A}(\zeta)(x)\in A_i\},\\
 \operatorname{Ext}_{i,\mathfrak A}(\zeta=\eta)
   &=\{x:\operatorname{Ext}_{i,\mathfrak A}(\zeta)(x)
          =\operatorname{Ext}_{i,\mathfrak A}(\eta)(x)\},\\
 \operatorname{Ext}_{i,\mathfrak A}(P[\bar\zeta])
   &=\{x:(\operatorname{Ext}_{i,\mathfrak A}(\zeta_r)(x))_{r<n}
                   \in F_P(i)\},\\
 \operatorname{Ext}_{i,\mathfrak A}(\neg\phi)
   &=U^\omega\setminus\operatorname{Ext}_{i,\mathfrak A}(\phi),\\
 \operatorname{Ext}_{i,\mathfrak A}(\phi\land\psi)
   &=\operatorname{Ext}_{i,\mathfrak A}(\phi)
                    \cap\operatorname{Ext}_{i,\mathfrak A}(\psi),\\
 \operatorname{Ext}_{i,\mathfrak A}(\exists v_n\phi)
   &=\{x:\text{some }y\in U\text{ has }
          x[n:=y]\in\operatorname{Ext}_{i,\mathfrak A}(\phi)\},\\
 \operatorname{Ext}_{i,\mathfrak A}(N\phi)
   &=\{x:(i,\{j\in I:x\in\operatorname{Ext}_{j,\mathfrak A}(\phi)\})
                  \in R_N\}.
\end{align*}
Disjunction is union, implication and biconditional have their Boolean
set operations, and universal quantification replaces ``some'' by ``every.''
All displayed sets are restricted to \(x\in U^\omega\).

A sentence is true at \(i\) iff its extension is \(U^\omega\).
If the only free variable of \(\phi\) is \(v_n\), then
\(y\operatorname{sat}_{i,\mathfrak A}\phi\) iff some assignment in
its extension has coordinate \(n\) equal to \(y\).
\(\operatorname{Int}_{\mathfrak A}(\xi)\) is the function
\(i\mapsto\operatorname{Ext}_{i,\mathfrak A}(\xi)\), for a variable,
individual constant, or formula. Sentence consequence is \emph{local}:
in every interpretation and at every one index, truth of all premises
implies truth of the conclusion at that same index. Logical truth is
consequence of no premises, and equivalence is validity of the biconditional.
\end{hbdefinition}

\begin{hblemma}{thm:p08-montague-pragmatic-recursion}{proved=annex:p8.n-b-006.2.pragmatic}{Well-defined recursion and the eight pragmatic criteria}
The simultaneous extensions at all indices are well defined. They depend
on assignment coordinates only for free variables; sentence extensions
are therefore \(U^\omega\) or \(\varnothing\).
For nonempty interpretations the definitions give the source's criteria:
\begin{enumerate}
\item \(P[c,d]\) is true iff \((F_c(i),F_d(i))\in F_P(i)\).
\item \(y\) satisfies \(P[c,v_n]\) iff \((F_c(i),y)\in F_P(i)\).
\item \(y\) satisfies \(c=v_n\) iff \(F_c(i)=y\).
\item \(y\) satisfies \(E[v_n]\) iff \(y\in A_i\).
\item Truth of \(\neg\phi\) is falsity of the sentence \(\phi\).
\item Truth of \(\phi\land\psi\) is truth of both sentences.
\item \(\exists v_n\phi\), for \(\phi\) with that one free variable,
      is true iff some possible individual satisfies \(\phi\).
\item \(N\phi\), for a sentence \(\phi\), is true iff
      \((i,\{j:\phi\text{ is true at }j\})\in R_N\).
\end{enumerate}
\end{hblemma}

Actualist quantification is expressible as
\(\exists v_n(E[v_n]\land\phi)\) or
\(\forall v_n(E[v_n]\to\phi)\); the primitive quantifier itself ranges
over \(U\). This permits reference to someone who existed earlier but
is no longer remembered. The extension of a closed individual constant
is the constant assignment-function with its denotation as value, and
the extension of a sentence is the assignment-set encoding of its truth
value. These encode Frege's ordinary extensions. Under an operator,
the whole intension of the operand is needed. There is no further
hierarchy of indirect intensions: one may instead recursively construct
intensions first, each as a function of its components' intensions
\cite[pp.~75--76]{Montague1970Pragmatics}.

\begin{hbtheorem}{thm:p08-montague-operator-examples}{proved=annex:p8.n-b-006.2.examples}{Past tense, relational modalities, and the progressive example}
In a nonempty interpretation on \(I=\mathbb R\), let \(R_N(i,J)\) mean
\(\exists j\in J\,(j<i)\). The operator then expresses past truth.
Replacing \(<\) by \(>\) gives future truth. Relational necessity
with accessibility \(r\) is obtained by requiring \(r[i]\subseteq J\).

Scott's example in the article instead takes \(R_N(i,J)\) to mean that
\(J\) includes an open interval containing \(i\). Then \(N\phi\)
holds exactly when \(\phi\) holds throughout some such interval. The
source offers this as a model for a progressive reading; it does not
claim a complete linguistic analysis of the English progressive. This
family cannot be represented by a fixed accessibility-successor set
with either the universal or existential relational modal clause.
\end{hbtheorem}

\subsection{Intensional logic: Definitions VIII--XIII}

\begin{hbdefinition}{def:p08-montague-intensional-language}{Types, predicate variables, and descriptions}
An intensional language retains individual constants, variables, \(E\),
identity, and connectives. It adds predicate variables \(G_{n,k}\) for
every arity \(k\geq0\), predicate constants of type
\(\langle s_0,\ldots,s_{r-1}\rangle\), with each \(s_q\geq-1\),
quantifiers for individual and predicate variables, the primitive
\(\Box\), and descriptions \(\mathsf T G\phi\).
Type \(-1\) denotes an individual argument; type \(k\geq0\) denotes
an argument that is a \(k\)-place predicate \emph{intension}.
An ordinary \(r\)-place predicate constant has all its argument types
equal to \(-1\). In particular a constant of type \(\langle0\rangle\)
is a unary predicate of propositions, not an ordinary zero-place constant.

Atoms are \(E[\zeta]\), \(\zeta=\eta\),
\(G[\zeta_0,\ldots,\zeta_{k-1}]\), and well-typed
\(P[\xi_0,\ldots,\xi_{r-1}]\). In the last expression an argument of
type \(-1\) is an individual symbol and an argument of type \(k\)
is a \(k\)-place predicate variable or description of that sort.
Formulas close under connectives, either sort of quantifier, and \(\Box\).
The source's final formation operation replaces occurrences of a predicate
variable in argument or predicate positions by \(\mathsf T G\phi\),
not occurrences immediately following a quantifier or description binder.
Thus descriptions can serve both as predicate terms and as higher-typed
arguments. Only predicate variables may be described; individual
descriptions and operation symbols are not part of this system.
\end{hbdefinition}

\begin{hbdefinition}{def:p08-montague-intensional-interpretation}{Full predicate domains and typed assignments}
A possible interpretation is \(\mathfrak A=\langle A,F\rangle\),
with the same \(I,A_i,U\) and individual-constant functions as before.
Put
\[
 \mathcal P_k=\{Y:Y\text{ is a function }I\to\mathcal P(U^k)\}.
\]
For a constant of type \(\langle s_0,\ldots,s_{r-1}\rangle\),
\(F_P(i)\subseteq\prod_{q<r}D_{s_q}\), where
\(D_{-1}=U\) and \(D_k=\mathcal P_k\). These are full predicate
domains, not Henkin-selected subfamilies.

An assignment system \(x\) has coordinates \((n,k)\), with
\(n\in\mathbb N\) and \(k\geq-1\), taking values in \(D_k\).
Let \(S\) be the set of all such systems; \(v_n\) uses coordinate
\((n,-1)\), and \(G_{n,k}\) uses \((n,k)\).
For a predicate term write \(\llbracket\xi\rrbracket_i^x\) for
\(\operatorname{Ext}_{i,\mathfrak A}(\xi)(x)\); it is a member of
\(\mathcal P_k\), not merely its extension at \(i\).
The empty predicate \(O_k\) has \(O_k(j)=\Lambda\) at every index;
its graph is the source's \(I\times\{\Lambda\}\).
\end{hbdefinition}

\begin{hbdefinition}{def:p08-montague-intensional-extension}{The twelve clauses of the typed extension recursion}
Extensions of terms are functions on \(S\), and extensions of formulas
are subsets of \(S\). The clauses are:
\begin{enumerate}
\item \(\llbracket c\rrbracket_i^x=F_c(i)\).
\item \(\llbracket v_n\rrbracket_i^x=x_{n,-1}\).
\item \(\llbracket G_{n,k}\rrbracket_i^x=x_{n,k}\).
\item Let \(B=\{Y\in\mathcal P_k:x[(n,k):=Y]\in
      \operatorname{Ext}_{i,\mathfrak A}(\phi)\}\).
      \(\llbracket\mathsf T G_{n,k}\phi\rrbracket_i^x\) is the sole
      member of \(B\) if there is exactly one, and \(O_k\) otherwise.
\item \(\operatorname{Ext}_{i,\mathfrak A}(E[\zeta])
      =\{x:\llbracket\zeta\rrbracket_i^x\in A_i\}\).
\item \(\operatorname{Ext}_{i,\mathfrak A}(\zeta=\eta)
      =\{x:\llbracket\zeta\rrbracket_i^x=\llbracket\eta\rrbracket_i^x\}\).
\item For a predicate variable or description \(\xi\) of arity \(k\),
      \(\operatorname{Ext}_{i,\mathfrak A}(\xi[\bar\zeta])
      =\{x:(\llbracket\zeta_q\rrbracket_i^x)_{q<k}
                         \in\llbracket\xi\rrbracket_i^x(i)\}\).
\item For a well-typed constant application,
      \(\operatorname{Ext}_{i,\mathfrak A}(P[\bar\xi])
      =\{x:(\llbracket\xi_q\rrbracket_i^x)_{q<r}\in F_P(i)\}\).
\item Negation is complement in \(S\), conjunction is intersection,
      and the other connectives have their Boolean operations.
\item \(\operatorname{Ext}_{i,\mathfrak A}(\exists v_n\phi)
      =\{x:\exists y\in U\ x[(n,-1):=y]\in
                      \operatorname{Ext}_{i,\mathfrak A}(\phi)\}\);
      universal quantification uses every \(y\in U\).
\item \(\operatorname{Ext}_{i,\mathfrak A}(\exists G_{n,k}\phi)
      =\{x:\exists Y\in\mathcal P_k\ x[(n,k):=Y]\in
                      \operatorname{Ext}_{i,\mathfrak A}(\phi)\}\);
      universal quantification uses every such \(Y\).
\item \(\operatorname{Ext}_{i,\mathfrak A}(\Box\phi)
      =\{x:\forall j\in I\ x\in\operatorname{Ext}_{j,\mathfrak A}(\phi)\}\).
\end{enumerate}
All sets here range over \(x\in S\). A sentence is true iff its extension
is \(S\). Single-variable satisfaction uses existence of a system in the
extension with the indicated individual or predicate coordinate. Intension
is again the function of the index giving these extensions. Consequence,
truth, and equivalence are defined locally as before, now over all
intensional interpretations and languages.
\end{hbdefinition}

\begin{hblemma}{thm:p08-montague-intensional-recursion}{proved=annex:p8.n-b-006.2.intensional}{The typed recursion and the additional criteria}
The simultaneous term/formula recursion is well defined and respects
free variables. For nonempty interpretations, besides the earlier seven
nonmodal criteria it gives:
predicate quantification over all \(\mathcal P_k\); application of a
type-\(\langle k\rangle\) constant to the whole predicate intension;
truth of \(\Box\phi\) at every index; unique-predicate description
with the empty-predicate default; and, for a proposition variable,
\(G[\ ]\) true iff its assigned proposition takes value \(\mathrm T\)
at that index. The primitive \(\Box\) is this universal-index modality,
not an arbitrary pragmatic operator \(N\).
\end{hblemma}

\begin{hbtheorem}{thm:p08-montague-concept-relations}{proved=annex:p8.n-b-006.2.concepts}{Individual concepts and extensional relations inside intensional logic}
In a nonempty interpretation, at an index \(i\in I\), functions
\(I\to U\) correspond exactly to unary predicates satisfying
\[
 \Box\exists u\forall v\,(G[v]\leftrightarrow v=u).
\]
Binary relations on \(U\) correspond exactly to binary predicates
satisfying
\[
 \forall u\forall v\,(\Box G[u,v]\lor\Box\neg G[u,v]).
\]
The same constant-intension construction applies to every arity, including
zero. These are definable restrictions of the predicate domains, not
additional primitive variable sorts.
\end{hbtheorem}

\begin{hbdefinition}{def:p08-montague-predicate-names}{Predicates expressed by formulas}
For distinct individual variables \(\bar u=(u_0,\ldots,u_{m-1})\)
and a fresh \(m\)-place predicate variable \(G\), define the source's
predicate name by
\[
 \widehat{\bar u}\phi
   =\mathsf T G\,\forall\bar u\,\Box(G[\bar u]\leftrightarrow\phi).
\]
Freshness means that \(G\) does not occur in \(\phi\); the source
chooses the first available variable of its sort. For \(m=0\),
\[
 \widehat\phi=\mathsf T G\,\Box(G[\ ]\leftrightarrow\phi)
\]
names the proposition that \(\phi\). An occurrence of \(\widehat\phi\)
is an abbreviation involving \(\Box\) and a description, not another
primitive logical constant.
\end{hbdefinition}

\begin{hbtheorem}{thm:p08-montague-descriptions}{proved=annex:p8.n-b-006.2.descriptions}{Description elimination and the interpretation of predicate names}
Every sentence is equivalent to one without descriptions. At a fixed
assignment, \(\widehat{\bar u}\phi\) denotes the unique predicate
\(Y\) with
\[
 Y(j)=\{\bar a\in U^m:x[\bar u:=\bar a]\in
                          \operatorname{Ext}_{j,\mathfrak A}(\phi)\}.
\]
In particular \(\widehat\phi\) denotes its truth function across all
indices, with the remaining assignment held fixed. With identity of
propositions available as an abbreviation, \(\Box\phi\) can be
recovered by equating \(\widehat\phi\) to
\(\widehat{\forall v_0(v_0=v_0)}\).
\end{hbtheorem}

\subsection{Operators as predicates of propositions}

\begin{hbtheorem}{thm:p08-montague-pragmatic-reduction}{proved=annex:p8.n-b-006.2.reduction}{The reduction of ordinary pragmatics to intensional logic}
For each pragmatic operator \(N\), add a distinct constant \(N'\)
of type \(\langle0\rangle\). Extend the original \(F\) by
\[
 F'_{N'}(i)=\{\langle Y\rangle:Y\in\mathcal P_0,
                (i,\{j:Y(j)=\mathrm T\})\in R_N\}.
\]
Preserve individual constants and ordinary predicate constants. Translate
connectives, quantifiers, and atoms compositionally, and translate
\(N\psi\) as
\[
 \exists G\,(\Box(G[\ ]\leftrightarrow\psi^*)\land N'[G]),
\]
with fresh zero-place \(G\). Then each sentence has the same truth at
each index before and after translation. This is the source's precise
reduction, not a translation by an accessibility relation.
\end{hbtheorem}

\begin{hbtheorem}{thm:p08-montague-converse-reduction}{proved=annex:p8.n-b-006.2.converse}{The converse correspondence and its syntactic restriction}
For a type-\(\langle0\rangle\) predicate constant \(P\), define
\[
 R_{P'}=\{(i,J):\exists Y\ (\langle Y\rangle\in F_P(i)
                         \land J=\{j:Y(j)=\mathrm T\})\}.
\]
Every predicate of propositions thereby determines a unary pragmatic
modality, and the two constructions are inverse because propositions
are determined by their truth sets.

For an intensional language whose predicate constants have either type
\(\langle0\rangle\) or only \(-1\) arguments, add pragmatic
operators \(P'\) and a further operator \(N\) with
\(R_N(i,J)\) iff \(J=I\). Replace \(P[\widehat\psi]\) by
\(P'\psi^*\), and \(\Box\psi\) by \(N\psi^*\), recursively.
Truth is preserved \emph{provided the resulting expression is a
pragmatic sentence}. This qualification excludes unrestricted predicate
quantification or predicate arguments that have not disappeared under
the indicated replacements. The article claims only this partial reverse
reduction, not equivalence of the full languages.
\end{hbtheorem}

In the terminology of this article, modal logic is the part of the
intensional language without predicate variables; intensional logic is
accordingly a second-order modal logic, and pragmatics a first-order
reduction of a part of it. Montague mentions further higher-order and
even transfinite extensions, but develops only the second-order system
needed for the applications here. This is a source boundary, not a
definition identifying all modern modal logics with this universal-index
fragment \cite[p.~84]{Montague1970Pragmatics}.

\begin{hbexample}{def:p08-montague-belief-seems}{Belief, iteration, and properties attributed to objects}
Let \(B\) have type \(\langle-1,0\rangle\), with
\((a,Y)\in F_B(i)\) interpreted as belief by \(a\) in proposition
\(Y\) at world \(i\). With individual constants \(J,R\) and an
ordinary unary \(S\), the source's iterated, quantified belief example is
\[
 \exists x\,(E[x]\land B[J,\widehat{B[R,\widehat{S[x]}]}]).
\]
A description-free equivalent is
\[
 \exists x\exists G\bigl(E[x]\land B[J,G]\land
 \Box(G[\ ]\leftrightarrow
       \exists H(B[R,H]\land\Box(H[\ ]\leftrightarrow S[x])))\bigr).
\]
For ``seems,'' a constant \(\mathit{Se}\) has type
\(\langle-1,1,-1\rangle\), and
\(\mathit{Se}[u,\widehat w S[w],v]\) relates two individuals to a
property intension.

These constants receive empirical interpretations, not logical
definitions of belief or seeming. The source offers assent to a sentence,
or assent to an open formula while its object is indicated, as fallible
confirmatory criteria; neither is proposed as conclusive or as a complete
behavioristic definition. If two sentences are logically equivalent,
their proposition names coincide, and belief in them is equivalent in
this semantics. Montague identifies the resulting objection and suggests
distinguishing another notion whose objects are sentences or complexes
containing open formulas. He does not provide that alternative system
here \cite[pp.~84--86; notes~10--11]{Montague1970Pragmatics}.
\end{hbexample}

\subsection{Extended pragmatics}

\begin{hbdefinition}{def:p08-montague-extended-pragmatics}{Degrees, bound variables, and the extended operator clause}
An extended pragmatic operator has degree \(\langle m,n,p\rangle\)
and forms expressions
\[
 N u_0\cdots u_{m-1}
       [\zeta_0,\ldots,\zeta_{n-1},\phi_0,\ldots,\phi_{p-1}],
\]
where the \(u_r\) are distinct individual variables bound in the
formula arguments, the \(\zeta_r\) are individual symbols, and
the \(\phi_q\) are formulas. Other formation operations are identity,
actuality, connectives, and individual quantifiers. An interpretation
is \(\langle A,F\rangle\), with
\[
 F_N(i)\subseteq U^n\times\mathcal P_m^{\,p}.
\]
The ordinary assignment space is again \(U^\omega\). If \(u_r=v_{k_r}\)
and \(x\in U^\omega\), put
\[
 Y_{q,x}(j)=\{(a_0,\ldots,a_{m-1})\in U^m:
       x[k_0:=a_0,\ldots,k_{m-1}:=a_{m-1}]
              \in\operatorname{Ext}_{j,\mathfrak A}(\phi_q)\}.
\]
The extended operator expression has extension consisting of precisely
the assignments \(x\) such that
\[
 (\operatorname{Ext}_{i,\mathfrak A}(\zeta_0)(x),\ldots,
  \operatorname{Ext}_{i,\mathfrak A}(\zeta_{n-1})(x),
  Y_{0,x},\ldots,Y_{p-1,x})\in F_N(i).
\]
The remaining extension clauses are the earlier actuality, identity,
Boolean, and individual-quantifier clauses. Truth and local consequence
use these extensions as in ordinary pragmatics.
\end{hbdefinition}

\begin{hbtheorem}{thm:p08-montague-extended-reduction}{proved=annex:p8.n-b-006.2.extended}{The two inclusions and their limits}
An operator of degree \(\langle0,n,0\rangle\) is an ordinary
\(n\)-place predicate; degree \(\langle0,0,1\rangle\) accommodates
an ordinary pragmatic operator by identifying propositions with truth
sets. An extended operator of arbitrary degree is represented in
intensional logic by a predicate constant with \(n\) arguments of
type \(-1\) followed by \(p\) arguments of type \(m\). Replace its
formula arguments by the corresponding names
\(\widehat{\bar u}\phi_q\). This translation preserves truth.

In particular belief becomes a degree-\(\langle0,1,1\rangle\)
operator \(B'[x,\phi]\), and seeming becomes a degree-
\(\langle1,2,1\rangle\) operator
\(\mathit{Se}'w[u,v,\phi]\). These reductions do not provide arbitrary
predicate variables or predicate quantifiers in extended pragmatics.
\end{hbtheorem}

The article further asserts that the unrestricted intensional sentence
that Jones believes something Robinson does not believe has no
corresponding extended-pragmatic formula. It does not give a formal
inexpressibility proof; that assertion is not being assigned the proved
status of the preceding constructive inclusion
\cite[p.~88]{Montague1970Pragmatics}.

\subsection{Metatheory and referential distinctions}

\begin{hbdefinition}{def:p08-montague-predicative}{The source's predicative fragment and effectiveness convention}
A predicative intensional sentence contains no descriptions. Every
predicate-universal subformula is of the form
\(\forall G(P[\bar\xi]\to\chi)\), and every predicate-existential
subformula is of the form \(\exists G(P[\bar\xi]\land\chi)\), where
\(P\) is a predicate constant, the arguments are individual symbols
or predicate variables of the appropriate types, and \(G\) occurs
among these arguments. This is a syntactic guard by a nonlogical
constant; it is not unrestricted full predicate quantification.

A set of sentences is satisfiable if all are true at one index in one
nonempty interpretation. Source languages need not be countable.
Effectiveness is asserted only inside a fixed countable collection of
expressions with a suitable G\"odel numbering, containing the logical
symbols, all variables, infinitely many constants of each required
arity/type and operators of each degree, and closed under concatenation.
A recursive language has a recursive symbol set in this collection.
\end{hbdefinition}

\begin{hbtheorem}{thm:p08-montague-metatheory-source}{cited=Montague1970Pragmatics}{The article's compactness and effectiveness assertions}
Full intensional logic is not compact, and some recursive intensional
languages have non-recursively-enumerable sets of logical truths.
Predicative intensional sentences are compact; in any recursive language
their logical truths are recursively enumerable. Ordinary and extended
pragmatics are compact and, for recursive languages, their logical
truths are recursively enumerable.

The article sketches these claims through second-order logic and
Henkin's \(\omega\)-order completeness theorem. It invokes Craig's
axiomatization theorem for recursive axiomatizations under detachment,
reports Kaplan's axiomatization of ordinary pragmatics, and reports
the natural axiomatization of predicative intensional logic as then open.
It gives neither complete derivations of these results nor the proposed
axiom systems. These are explicit source proof boundaries, not modern
neighborhood-completeness theorems supplied by this entry
\cite[pp.~88--90; note~13]{Montague1970Pragmatics}.
Its note~13 further asserts the predicative compactness result for a
higher-order modal language with variables at all finite levels; it does
not provide a complete higher-order development in this article.
\end{hbtheorem}

\begin{hbtheorem}{thm:p08-montague-substitution}{proved=annex:p8.n-b-006.2.substitution}{Rigid-variable replacement and nonrigid-constant failure}
In all three systems, replacing a free occurrence of an individual variable
\(u\) by an individual variable \(v\), without capture, preserves
truth when their assigned values agree. Thus
\(\forall u\forall v(u=v\to(\phi\leftrightarrow\phi'))\) is valid.
For individual constants \(c,d\), the corresponding
\(c=d\to(\phi\leftrightarrow\phi')\) is not generally valid.
Universal instantiation by a variable is valid, but instantiation by a
nonrigid individual constant is not generally valid.
\end{hbtheorem}

Montague applies this distinction to Morning Star/Evening Star belief
contexts. His historical discussion contrasts Church's system, Kaplan's
revision using Carnap's model-based worlds, and later Howard/Kaplan/Scott
systems. The source criticizes their restrictions on unrestricted
individual quantification and, in the earlier framework, belief
iteration. These are claims of the 1970 article, not blanket claims
about later systems. The paper reports a January 1967 talk, an earlier
submission withdrawn after a lengthy delay, and an extended-pragmatics
addition after discussions of Scott's June 1967 treatment
\cite[pp.~90--94]{Montague1970Pragmatics}.

\subsection{The justified bridge to modern neighborhoods}

\begin{hbtheorem}{thm:p08-montague-neighborhood-bridge}{proved=annex:p8.n-b-006.2.neighborhood}{From Montague's relation to a neighborhood function}
For a fixed pragmatic operator in a nonempty interpretation with a
nonempty index set, define
\[
 \mathcal N_N(i)=\{J\subseteq I:(i,J)\in R_N\}.
\]
Then its sentence clause is exactly the modern neighborhood clause:
\(N\phi\) holds at \(i\) iff its truth set belongs to
\(\mathcal N_N(i)\). Conversely every neighborhood function determines
such an \(R_N\). On the propositional fragment this gives a
truth-preserving correspondence of models, not merely a similarity
of terminology. No upward, intersection, or unit closure is required
by Montague's possible-interpretation definition.
\end{hbtheorem}

In this comparison M abbreviates
\(\Box(p\land q)\to(\Box p\land\Box q)\), C abbreviates
\((\Box p\land\Box q)\to\Box(p\land q)\), and N abbreviates
\(\Box\top\). The congruential base E uses classical propositional
reasoning and replacement of theorem-level equivalents under the box.
The existing exact calculus sheets, rather than new rule displays here,
fix these presentations:
\cref{sheet:hilbert-e,sheet:hilbert-m,sheet:hilbert-c,sheet:hilbert-n,sheet:hilbert-r,sheet:hilbert-k}.

\begin{hbtheorem}{thm:p08-montague-modern-boundary}{proved=annex:p8.n-b-006.2.modern}{Closure and consequence distinctions in the modern comparison}
In the modern propositional notation, upward closure gives M,
binary-intersection closure gives C, and containing the whole index set
gives N. Replacement of globally equivalent formulas follows already
from dependence on truth sets. The combination EMC is regular but
need not contain the unit; EN need not be monotone or intersection
closed. The normal combination is EMCN, equivalent at theoremhood
to K with its normal necessitation rule. EN alone is not K.

These names belong to the modern comparison, not to the syntax or
axiomatization of the 1970 paper. Its logical consequence is local;
necessitation does not preserve arbitrary local premises. The source's
primitive universal-index \(\Box\) and its arbitrary operators \(N\)
must likewise remain distinct.
\end{hbtheorem}

For modern conventions and the separately allocated completeness
constructions see the neighborhood presentation in
\cref{def:p08-neighborhood-model,def:p08-neighborhood-conditions}
\cite{Chellas1980,Pacuit2017}. The generalized relation itself does not
establish any of those completeness theorems, and the joint history
does not identify Montague's individual quantification with Scott's
quantification over individual concepts.
